% Compile with: uplatex file.tex && dvipdfmx file.dvi
%           or: latexmk -uplatex file.tex
\documentclass[dvipdfmx]{article}
\usepackage{amsmath, amssymb}
\usepackage{tikz}

\title{scenario\_rlc\_transfer\_function}
\date{\today}

\begin{document}
\maketitle

Consider a series RLC circuit (charge $q(t)$, input voltage $v(t)$). Prepare $q,v$ as time-dependent functions.

Define q as a time-dependent function

$\checkmark$ Declared q as time-dependent (time variable: t)

Define v as a time-dependent function

$\checkmark$ Declared v as time-dependent (time variable: t)

Name the Laplace transform of $q$ as $Q$ (an s-domain quantity), via a definition card.

Let Q be the Laplace transform of q

\begin{equation}
  Q(s) = \int_0^{\infty} q(t)\, e^{-s t}\, dt
  \label{eq:E1}
\end{equation}

Likewise name the Laplace transform of $v$ as $V$.

Let V be the Laplace transform of v

\begin{equation}
  V(s) = \int_0^{\infty} v(t)\, e^{-s t}\, dt
  \label{eq:E2}
\end{equation}

Write the Laplace transform rules for the derivatives. Integration by parts here fixes $u=e^{-st}$ by the structure of the transform, so there is no choice to make: this is a theorem, and the researcher writes it. First $\mathcal L\{\dot q\} = sQ - q(0)$.

$s Q - q_0$

\begin{equation}
  s Q - q_0
  \label{eq:E3}
\end{equation}

Then $\mathcal L\{\ddot q\} = s^2 Q - s q(0) - \dot q(0)$.

$s^2 Q - s q_0 - u_0$

\begin{equation}
  s^2 Q - s q_0 - u_0
  \label{eq:E4}
\end{equation}

Assemble the s-domain form of the circuit equation $L\ddot q+R\dot q+q/C=v$ (inductance $L$, resistance $R$, capacitance $C$), referring to the two rules above by their E-ids.

$L E4 + R E3 + Q/C = V$

\begin{equation}
  L \left(-u_{0}+Q s^2-q_{0} s\right)+R \left(Q s-q_{0}\right)+\frac{Q}{C}=V
  \label{eq:E5}
\end{equation}

The transfer function is defined under zero initial conditions, so state $q(0)=0$ and $\dot q(0)=0$. An H-1 definition is sent to the CAS when registered, so these are in force from here on.

$q_0 = 0$

\begin{equation}
  q_0 = 0
  \label{eq:E6}
\end{equation}

$u_0 = 0$

\begin{equation}
  u_0 = 0
  \label{eq:E7}
\end{equation}

Define the transfer function $H=Q/V$ (the s-domain ratio of output to input).

$H = Q/V$

\begin{equation}
  H = Q/V
  \label{eq:E8}
\end{equation}

Solve the s-domain equation for $Q$. The zero initial conditions are already in force.

Solve E5 for Q

\begin{equation}
  Q = \frac{C V}{C L s^2+C R s+1}
  \label{eq:E9}
\end{equation}

Substitute the resulting expression for $Q$ into the definition of $H$ to get the closed form of the transfer function.

Substitute E9 into E8

\begin{equation}
  H = \frac{C}{C L s^2+C R s+1}
  \label{eq:E10}
\end{equation}

\end{document}